\section{境界外のsheafと拡大によるcohomologyの相殺}
\label{sec:offboundary-extensions}
本節は既約根sheafに限らない．既知入力は
Rainsの境界triangle \eqref{eq:boundarytriangle}，
$1$次元以下のsheafへの$\Ext^2$消滅
\cite[Corollary 7.23]{RainsBirational}，および
中心正則元$g$を持つ反標準環の仮定である．
$\mathcal Z$をRains heartの中で
\[
 \mathcal Z=\{M\in\mathcal A_{\mathrm R}:
    \dim M\le1,\ Li^*M\simeq0\}
\]
と置く．$R$は右noetherianであり，必要な$\chi$条件により
$\Phi=\pi\mathbf R\Gamma_*$を該当対象に定義できると仮定する．
$A=R/(g-1)\cong (R[g^{-1}])_0$と書く．
ここで$\pi$は有限次元次数付き加群によるSerre商である．

\begin{theorem}[境界外のsheafに対する二項振幅と長完全列]
\label{thm:offboundary-amplitude}
$\mathcal Z$は拡大に閉じたexact categoryである．
任意の$M\in\mathcal Z$について次が成立する：
\begin{enumerate}
\item 自然な$\theta M\to M$は同型である．従って
$\tau^nM\simeq M$で，全$n\in\mathbb Z$について
$\Ext^j(e,\tau^n M)=0$ $(j\notin\{0,1\})$．
\item
\[
 N^j(M)=\bigoplus_{n\ge0}\Ext^j(e,\tau^nM)
 \qquad(j=0,1)
\]
は次数付き右$R$加群で，$g$による乗法は全次数で同型である．
特に$k[g]$加群として有限生成freeで，そのrankは
$h^j(M)=\dim_k\Ext^j(e,M)$．
\[
 V^j(M)=(N^j(M)[g^{-1}])_0\in\operatorname{mod}_{fd}A,
 \qquad \dim_kV^j(M)=h^j(M)
\]
であり，$H^j(\Phi M)\simeq\pi N^j(M)$．
従って自然な導来切断のcohomologyは次数$0,1$だけで，
どちらも反標準境界$g=0$の外にある．
\item 任意の短完全列$0\to M_1\to M_2\to M_3\to0$
（$M_i\in\mathcal Z$）から，有限次元右$A$加群の
自然な完全列
\begin{equation}\label{eq:offboundary-six}
0\to V^0(M_1)\to V^0(M_2)\to V^0(M_3)
\xrightarrow{\delta} V^1(M_1)
\to V^1(M_2)\to V^1(M_3)\to0
\end{equation}
を得る．この$\delta$はextension classのconnecting morphismで，
$A$-線型である．
\end{enumerate}
\end{theorem}
\begin{proof}
Rainsのdimension filtrationは短完全列に関し安定である．
また$Li^*$はexact triangulated functorなので
両端の制限が零なら中間項の制限も零となる．
従って$\mathcal Z$は拡大に閉じる．

境界triangle \eqref{eq:boundarytriangle}と$Li^*M=0$から
自然な$\theta M\to M$が同型となる．
$\theta$，$\tau$はheart上でexactなので，全twistも
$1$次元以下のsheafとなる．
負のExtは二つのheart対象間で消え，Rainsの
\cite[Corollary 7.23]{RainsBirational}で$\Ext^2$が消える．
この曲面のcohomological dimensionは高々$2$だから，
(1)が従う．

境界triangleの$\theta\to\mathrm{id}$から得られる
次数$1$の中心元$g$は$M$上で同型であり，
$\Ext^j(e,\tau^nM)\to\Ext^j(e,\tau^{n+1}M)$に
同型を誘導する．有限次元性により
$N^j(M)\simeq k[g]\otimes_k\Ext^j(e,M)$
が$k[g]$加群として成り立つ（$R$作用はこの分解で
単純になるとは主張しない）．
ゆえに$\pi N^j$は非零なら境界外の対象で，
dehomogenizationは$A$の有限次元加群$V^j$を与える．
$\Phi M$のcohomologyの同定は$\pi$のexactnessと
$\mathbf R\Gamma_*$の定義による．

短完全列の$\tau^n$-twistに$\RHom(e,-)$を適用すると
cohomologyは次数$0,1$だけなので六項完全列になる．
これを全$n\ge0$で直和し，中心元$g$で局所化した後
次数$0$を取る．局所化と次数$0$の抽出はこの
graded-moduleのクラスでexactであり，
自然性によりconnecting morphismは$A$-線型となる．
\end{proof}

\begin{corollary}[拡大に閉じた正・負の部分圏]
\label{cor:offboundary-sign-exact}
\[
 \mathcal Z_+=\{M\in\mathcal Z:\Ext^1(e,M)=0\},\qquad
 \mathcal Z_-=\{M\in\mathcal Z:\Hom(e,M)=0\}
\]
はともに拡大に閉じる．
$\Phi(\mathcal Z_+)$はanticanonical heartの次数$0$に，
$\Phi(\mathcal Z_-)$はその$[-1]$ shiftに含まれる．
$V^0:\mathcal Z_+\to\operatorname{mod}_{fd}A$，
$V^1:\mathcal Z_-\to\operatorname{mod}_{fd}A$
はそれぞれ短完全列を保つ．
特に正指数（負指数）の既約根sheafの任意の有限回の
拡大は$\mathcal Z_+$（$\mathcal Z_-$）に入る．
\end{corollary}
\begin{proof}
$\Ext^1$，$\Hom$の各々の消滅は
\eqref{eq:offboundary-six}で拡大に関して閉じる．
定理\ref{thm:sign}により各既約根sheafは指定部分圏に属する．
長完全列の端の項が零なら残る三項は短完全列となる．
\end{proof}

\begin{proposition}[正負の混合拡大と相殺]
\label{prop:mixed-sign-cancellation}
$m_+>0$，$m_-<0$を持つ境界外既約根sheaf
$M_+\in\mathcal Z_+$，$M_-\in\mathcal Z_-$が同じ
Rains heartに実在すると仮定する．
$0\to M_-\to M\to M_+\to0$という拡大を取る．
この拡大が誘導する$A$-線型写像
\[
 \delta_M:V^0(M_+)\longrightarrow V^1(M_-)
\]
のrankを$r$とすると，
\begin{equation}\label{eq:mixed-ranks}
 \dim_kV^0(M)=m_+-r,\qquad
 \dim_kV^1(M)=-m_--r .
\end{equation}
特に$m_+=-m_->0$の場合，$\delta_M$が同型なら
$\Phi(M)=0$である．
従って$\Phi$が$\mathcal Z$上で保守的なら，
非零の混合拡大に対するそのような同型
$\delta_M$は存在しない．
同じ結論は$\Phi$が全導来圏の同値であるなら必ず成立する．
\end{proposition}
\begin{proof}
\eqref{eq:offboundary-six}に
$V^0(M_-)=V^1(M_+)=0$を代入すると
\[
0\to V^0(M)\to V^0(M_+)
\xrightarrow{\delta_M}V^1(M_-)\to V^1(M)\to0
\]
となり，\eqref{eq:mixed-ranks}が従う．
$\delta_M$が同型なら$N^0(M),N^1(M)$は
次数$0$の局所化が零であるが，$g$は各次数において
同型なので両方零であり，$\Phi(M)=0$である．
$M$は$M_-$を非零部分対象とするため非零である．
\end{proof}

\begin{remark}[Euler標数だけでは検出できない]
一般の$M\in\mathcal Z$に対して
\[
 \chi(e,M)=\dim V^0(M)-\dim V^1(M)
\]
であるが，$\chi(e,M)=0$から$\Phi(M)=0$は従わない．
正負の拡大類によるconnecting morphism
$\delta_M$のrankを指定して初めて
反標準切断のcohomologyが決まる．
$\mathcal Z_+,\mathcal Z_-$が$\mathcal Z$全体の
torsion pairになるとは証明していない．
\end{remark}

\begin{proposition}[根系からの共鳴交叉制約]
\label{prop:root-intersection}
Rains heartに境界に制限して零となる
異なる二つの既約純一次元根sheaf$M_\alpha,M_\beta$が
存在し，第一Chern類を$\alpha\ne\beta$とする．
すると
\[
 \Hom(M_\alpha,M_\beta)=\Ext^2(M_\alpha,M_\beta)=0,
 \qquad
 \dim_k\Ext^1(M_\alpha,M_\beta)=\alpha\cdot\beta.
\]
従って$\alpha\cdot\beta<0$ならそのような二対象は
同じRains heartには同時に存在できず，
$\alpha\cdot\beta=0$なら双方の向きの$\Ext^1$が零である．
$\alpha\cdot\beta=1$なら$\Ext^1$は一次元である．
\end{proposition}
\begin{proof}
非零な$M_\alpha\to M_\beta$の像$I$は，
$M_\beta$の純一次元性から一次元である．
$M_\alpha$の既約性より，核が一次元なら
その第一Chern類は$\alpha$となり，
像は零次元となって矛盾する．
核が零次元なら$M_\alpha$の純一次元性から零である．
従って写像は単射で，像の第一Chern類は$\alpha$である．
他方，$M_\beta$の既約性から非零部分sheafの第一Chern類は
$\beta$である．ゆえに$\alpha=\beta$となり矛盾する．

$\theta M_\alpha\simeq M_\alpha$であるから，
Serre双対により
$\Ext^2(M_\alpha,M_\beta)
 \simeq\Hom(M_\beta,\theta M_\alpha)^*=0$．
Rains\cite[Corollary 7.13, (7.4)]{RainsBirational}の
rank-zero Euler形式は
$\chi(M_\alpha,M_\beta)=-\alpha\cdot\beta$であり，
残る$\Ext^1$の次元は$\alpha\cdot\beta$となる．
\end{proof}

\begin{corollary}[導来完全忠実性に対する有限次元表現テスト]
\label{cor:affine-ext-test}
定理\ref{thm:offboundary-amplitude}の自然な導来切断$\Phi$が
二つの境界外既約根sheaf$M_+,M_-$（指数$m_+>0$，
$m_-<0$）の間の$\Ext^1$で完全忠実ならば，
対応する有限次元$A$加群$V_+$，$V_-$について
\[
\Ext^1_X(M_+,M_-)
 \xrightarrow{\;\sim\;}\Hom_A(V_+,V_-)
\]
が成り立つ．$\alpha_+=c_1(M_+)\ne c_1(M_-)=\alpha_-$なら
\[
\boxed{\dim_k\Hom_A(V_+,V_-)=\alpha_+\cdot\alpha_- .}
\]
従って右辺の次元がこれと異なるなら，
$\Phi$はその二対象の間で完全忠実ではない．
\end{corollary}
\begin{proof}
$\Phi(M_+)\simeq V_+$，$\Phi(M_-)\simeq V_-[-1]$
である（$V_\pm$は境界外の有限支持対象として見る）．
したがって
\[
\Hom_{\Db(\qgr R)}(\Phi M_+,\Phi M_-[1])
 =\Hom_{\qgr R}(V_+,V_-)=\Hom_A(V_+,V_-).
\]
最後の同型は，$g$が各$V_\pm$上で可逆であり，
その次数付き表示が$k[g]$に関して有限自由であることから，
局所化がこの二対象間の$\Hom$を保存するためである．
$\Ext^1_X(M_+,M_-)$は左辺へ$\Phi$が誘導する空間であり，
完全忠実性を仮定すると同型となる．
命題\ref{prop:root-intersection}を使えば次元式を得る．
この写像は命題\ref{prop:mixed-sign-cancellation}の
$\delta_M$のextension classによる線型写像そのものである．
\end{proof}

\begin{proposition}[根の和による反対指数の同時共鳴排除]
\label{prop:root-sum-exclusion}
古典的avatar$S$がgenuine del Pezzoで，
滑らかな反標準楕円曲線$E$を持つとする．
二つの根$\alpha,\beta$に
\[
 \mathcal O_E(\alpha)\otimes q^{m_\alpha}\simeq\mathcal O_E,
 \qquad
 \mathcal O_E(\beta)\otimes q^{m_\beta}\simeq\mathcal O_E
\]
という共鳴があるとする．もし$\alpha\cdot\beta=1$なら，
$m_\alpha+m_\beta=0$は不可能である．
特に，隣接する$A_2$部分根系の根に対する
指数$+r$と$-r$の共鳴は同時には起こらない．
\end{proposition}
\begin{proof}
$\gamma=\alpha+\beta$は
$\gamma^2=-2-2+2=-2$，$H\cdot\gamma=0$を
満たすから再び根である．
$m_\alpha+m_\beta=0$なら二つの共鳴式の積から
$\mathcal O_E(\gamma)\simeq\mathcal O_E$となる．
しかしgenuine del Pezzoについての
命題\ref{prop:zero-index}は任意の根の
$E$への制限が非自明であることを示している．
従って矛盾する．
\end{proof}

\begin{remark}[完全な相殺を探すときの制限]
異なる根sheafの間の非自明な拡大は，
上の命題から$\alpha\cdot\beta>0$を必要とする．
特に最も単純な$\alpha\cdot\beta=1$の場合，
指数が互いに反対で次元も等しいような混合拡大は，
genuine classical avatarではそもそも同時共鳴できない．
これは$\delta_M$の同型による導来的消失が
genuineな族に現れないことの\emph{完全な証明ではない}．
別のroot対や非既約対象を排除できていないためである．
\end{remark}

\begin{remark}[古典的な$(-2)$曲線上の完全な相殺]
\label{ex:euler-cancellation}
滑らかな弱del Pezzo $S$に$(-2)$曲線
$i:D\simeq\mathbb P^1\hookrightarrow S$があり，
反標準曲線$E$は$D$と交わらないとする．
$H=-K_S$の$D$への制限は自明である．
$\mathbb P^1$のEuler完全列
\[
0\to\mathcal O_D(-2)\xrightarrow{(-y,x)}
 \mathcal O_D(-1)^{\oplus2}
 \xrightarrow{(x,y)}\mathcal O_D\to0
\]
を$i_*$で送る．左端を$M_-$，右端を$M_+$
とすると$\chi(M_-)=-1$，$\chi(M_+)=1$である．
中間項$M=i_*\mathcal O_D(-1)^{\oplus2}$は非零だが
\[
\RHom_S(\mathcal O_S(nH),M)
 \simeq\mathbf R\Gamma(D,\mathcal O_D(-1)^{\oplus2})=0
 \quad(n\in\mathbb Z).
\]
したがって$\delta_M:k\to k$は同型で，
反標準導来切断は$M$を零に送る．
これは前の根壁の導来情報消失を
「正負指数の拡大の完全な相殺」として説明する
古典的な模型である．\emph{この例はweak avatarの例であり，
genuine非可換曲面で同じ拡大が存在するという主張ではない．}
\end{remark}
